\newpage
\section{Line Search Methods}
Two concepts they mention in this chapter which I felt somewhat unfamiliar with were \text{Lipschitz} continuity and small $o$ notation.
Therefore, I thought I would introduce them both here before moving on to solving the exercises.
% TODO(martins): Cover the above concepts

\subsection{Solutions}
The selected exercises for this chapter are 3.2, 3.5, 3.7 and 3.8.

\subsubsection{Exercise 3.2}
The Wolfe conditions are the following two inequalities
\begin{equation}
    \begin{aligned}
        f(x_k + \alpha p_k)              & \leq f(x_k) + \alpha c_1 \nabla f(x_k)^T p_k \quad \textit{(Sufficient decrease condition)} \\
        \nabla f(x_k + \alpha p_k)^T p_k & \geq c_2 \nabla f(x_k)^T p_k  \quad \textit{(Curvature condition)}
    \end{aligned}
    \label{eq:wolfe_conditions}
\end{equation}
for $0 < c_1 < c_2 < 1$.
Let us find a situation where $c_2 < c_1$ implies no nonzero step length is possible.
We begin by first restating the Wolfe conditions in terms of $\phi(\alpha) := f(x_k + \alpha p_k)$

\begin{align*}
    \phi(\alpha)  & \leq \phi(0) + \alpha c_1 \phi'(0) \\
    \phi'(\alpha) & \geq c_2 \phi'(0)
\end{align*}

If we assume for example that $\phi(\alpha) = -\alpha + \alpha^2$, these conditions become
\begin{align*}
    -\alpha + \alpha^2   \leq -\alpha c_1 \implies \alpha & \leq 1 - c_1           \\
    -1 + 2\alpha  \geq -c_2 \implies \alpha               & \geq \frac{1 - c_2}{2}
\end{align*}

These imply the chained inequality
\[
    \frac{1 - c_2}{2} \leq \alpha \leq 1 - c_1
\]
These inequalities doesn't rule out all values of $c_2 < c_1$, take $c_2 = 0.5, c_1 = 0.6$ for example
\[
    0.25 \leq \alpha \leq 0.4
\]
but for a large difference, $c_2 = 0.1, c_1 = 0.9$,
\[
    0.45 \leq \alpha \leq 0.1
\]
no step size is admissible.


To generalize the proof for a broader class of convex functions, we parameterize a ``general'' convex function by defining its limits on curvature. Instead of looking at a specific function, we study a set of functions which are $m$-strongly convex and are $L$-smooth.
These properties mean that $\phi''$ has some lower bound and upper bounds. The lower bound is $m$ (m-strongly convex) and the upper bound is $L$ (L-smooth means the gradient is Lipschitz continuous):
\[
    0 < m \leq \phi''(\alpha) \leq L
\]
Since $p_k$ is a descent direction, we also know $\phi'(0) < 0$.
Because $\phi(\alpha)$ is $m$-strongly convex, we can bound the function from below using a Taylor expansion:
\[
    \phi(\alpha) \geq \phi(0) + \alpha \phi'(0) + \frac{m}{2}\alpha^2
\]
Replacing $\phi(\alpha)$ in the sufficient decrease condition with this lower bound,
\begin{align*}
    \phi(0) + \alpha \phi'(0) + \frac{m}{2}\alpha^2 & \leq \phi(0) + \alpha c_1 \phi'(0) \\
    \frac{m}{2}\alpha^2                             & \leq \alpha \phi'(0)(c_1 - 1)
\end{align*}
Since $\alpha > 0$, we can divide by $\alpha$:
\[
    \frac{m}{2}\alpha \leq \phi'(0)(c_1 - 1) \implies \alpha \leq \frac{2(c_1 - 1)\phi'(0)}{m}
\]
This is the maximum possible step size allowed by the sufficent decrease condition for this class of functions.
Because $\phi(\alpha)$ is $L$-smooth, its derivative cannot increase faster than $L$. We can bound the slope from above:
\[
    \phi'(\alpha) \leq \phi'(0) + L\alpha
\]
For the curvature condition to hold, the slope $\phi'(\alpha)$ must be greater than or equal to $c_2 \phi'(0)$. Therefore, our upper bound on the slope must also satisfy this:
\[
    \phi'(0) + L\alpha \geq c_2 \phi'(0)
\]
Rearranging for $\alpha$:
\[
    \alpha \geq \frac{(c_2 - 1)\phi'(0)}{L}
\]

This is the minimum possible step size required by the curvature condition.
An admissible step length $\alpha$ must satisfy both bounds simultaneously. Setting up the chained inequality:
\[
    \frac{(c_2 - 1)\phi'(0)}{L} \leq \alpha \leq \frac{2(c_1 - 1)\phi'(0)}{m}
\]
Removing the middle inequalities with $\alpha$ and dividing by $\phi'(0)$, we can see the relationship between $c_1$ and $c_2$:
\[
    \frac{c_2 - 1}{L} \geq \frac{2(c_1 - 1)}{m}
\]
Let's multiply both sides by $m/2$. In optimization, the ratio $\kappa = \frac{L}{m}$ is known as the condition number of the function ($\kappa \geq 1$). This gives us the admissibility rule:
\[
    c_1 - 1 \leq \frac{c_2 - 1}{2\kappa}
\]

\subsubsection{Exercise 3.5}
\subsubsection{Exercise 3.7}
\subsubsection{Exercise 3.8}